\documentclass[aspectratio=169,10pt]{beamer}

\usepackage{iftex}
\providecommand{\dunenoirsetupfonts}{}
\newcommand{\DuneNoirLocalFontPath}{fonts}

\newcommand{\DuneNoirMonoFallback}{%
  \IfFontExistsTF{JetBrainsMono Nerd Font Propo}{
    \setmonofont{JetBrainsMono Nerd Font Propo}[Scale=0.92]
  }{
    \IfFontExistsTF{JetBrainsMono Nerd Font}{
      \setmonofont{JetBrainsMono Nerd Font}[Scale=0.92]
    }{
      \setmonofont{Courier New}[Scale=0.92]
    }
  }
}

\newcommand{\DuneNoirSansFallback}{%
  \IfFontExistsTF{IBM Plex Sans}{
    \setsansfont{IBM Plex Sans}
    \setmainfont{IBM Plex Sans}
  }{
    \IfFontExistsTF{Helvetica Neue}{
      \setsansfont{Helvetica Neue}
      \setmainfont{Helvetica Neue}
    }{
      \setsansfont{Arial}
      \setmainfont{Arial}
    }
  }
}

\newcommand{\DuneNoirUseHelvetica}{%
  \renewcommand{\dunenoirsetupfonts}{%
    \IfFontExistsTF{Helvetica Neue}{
      \setsansfont{Helvetica Neue}
      \setmainfont{Helvetica Neue}
    }{
      \IfFontExistsTF{Helvetica}{
        \setsansfont{Helvetica}
        \setmainfont{Helvetica}
      }{
        \DuneNoirSansFallback
      }
    }
    \DuneNoirMonoFallback
  }%
}

\newcommand{\DuneNoirUseOpenSans}{%
  \renewcommand{\dunenoirsetupfonts}{%
    \IfFileExists{\DuneNoirLocalFontPath/OpenSans.ttf}{
      \setsansfont{Open Sans Custom}[
        Path=\DuneNoirLocalFontPath/,
        UprightFont={OpenSans.ttf},
        ItalicFont={OpenSans-Italic.ttf},
        BoldFont={OpenSans.ttf},
        BoldItalicFont={OpenSans-Italic.ttf},
        BoldFeatures={FakeBold=1.45},
        BoldItalicFeatures={FakeBold=1.45}
      ]
      \setmainfont{Open Sans Custom}[
        Path=\DuneNoirLocalFontPath/,
        UprightFont={OpenSans.ttf},
        ItalicFont={OpenSans-Italic.ttf},
        BoldFont={OpenSans.ttf},
        BoldItalicFont={OpenSans-Italic.ttf},
        BoldFeatures={FakeBold=1.45},
        BoldItalicFeatures={FakeBold=1.45}
      ]
    }{
      \DuneNoirSansFallback
    }
    \DuneNoirMonoFallback
  }%
}

\newcommand{\DuneNoirUseDosis}{%
  \renewcommand{\dunenoirsetupfonts}{%
    \IfFileExists{\DuneNoirLocalFontPath/Dosis.ttf}{
      \setsansfont{Dosis Custom}[
        Path=\DuneNoirLocalFontPath/,
        UprightFont={Dosis.ttf},
        ItalicFont={Dosis.ttf},
        BoldFont={Dosis.ttf},
        BoldItalicFont={Dosis.ttf},
        ItalicFeatures={FakeSlant=0.16},
        BoldFeatures={FakeBold=1.6},
        BoldItalicFeatures={FakeBold=1.6,FakeSlant=0.16}
      ]
      \setmainfont{Dosis Custom}[
        Path=\DuneNoirLocalFontPath/,
        UprightFont={Dosis.ttf},
        ItalicFont={Dosis.ttf},
        BoldFont={Dosis.ttf},
        BoldItalicFont={Dosis.ttf},
        ItalicFeatures={FakeSlant=0.16},
        BoldFeatures={FakeBold=1.6},
        BoldItalicFeatures={FakeBold=1.6,FakeSlant=0.16}
      ]
    }{
      \DuneNoirSansFallback
    }
    \DuneNoirMonoFallback
  }%
}

\newcommand{\DuneNoirUseKoHo}{%
  \renewcommand{\dunenoirsetupfonts}{%
    \IfFileExists{\DuneNoirLocalFontPath/KoHo-Regular.ttf}{
      \setsansfont{KoHo Custom}[
        Path=\DuneNoirLocalFontPath/,
        UprightFont={KoHo-Regular.ttf},
        BoldFont={KoHo-Bold.ttf},
        ItalicFont={KoHo-Italic.ttf},
        BoldItalicFont={KoHo-BoldItalic.ttf}
      ]
      \setmainfont{KoHo Custom}[
        Path=\DuneNoirLocalFontPath/,
        UprightFont={KoHo-Regular.ttf},
        BoldFont={KoHo-Bold.ttf},
        ItalicFont={KoHo-Italic.ttf},
        BoldItalicFont={KoHo-BoldItalic.ttf}
      ]
    }{
      \DuneNoirSansFallback
    }
    \DuneNoirMonoFallback
  }%
}

\DuneNoirUseHelvetica
\makeatletter
\input{../dune_dark_template_lab/beamerthemeDUNENoir.sty}
\makeatother

\ifPDFTeX
  \usepackage[T1]{fontenc}
  \usepackage[utf8]{inputenc}
\fi

\usepackage{amsmath,amssymb}
\usepackage{booktabs,array}
\usepackage{graphicx}
\usepackage{hyperref}
\usepackage{listings}
\usepackage{tikz}
\usetikzlibrary{arrows.meta,calc,positioning,fit,shapes.geometric}

\title[DAPHNE DAQ Integration]{DAPHNE DAQ Integration}
\subtitle{Ethernet data taking, standalone calibration, and \texttt{pds-run} scans}
\author{Manuel Arroyave}
\date{May 8, 2026}
\newcommand{\mydate}{May 8, 2026}

\lstdefinestyle{terminal}{
  basicstyle=\ttfamily\scriptsize,
  columns=fullflexible,
  keepspaces=true,
  breaklines=true,
  frame=single,
  rulecolor=\color{noirline},
  backgroundcolor=\color{noirpanel},
  xleftmargin=0.2em,
  xrightmargin=0.2em
}

\newcommand{\sectioncard}[2]{%
  \begin{frame}[plain]
    \begin{tikzpicture}[remember picture,overlay]
      \fill[noirpanelalt] (current page.south west) rectangle (current page.north east);
      \fill[duneorange] ($(current page.north west)+(0.9,-0.9)$) rectangle ++(1.8,0.08);
      \fill[noircyan] ($(current page.south east)+(-2.8,0.9)$) rectangle ++(1.8,0.08);
      \draw[noirline, line width=0.8pt] ($(current page.south west)+(0.9,1.1)$) -- ($(current page.south east)+(-0.9,1.1)$);
    \end{tikzpicture}
    \vspace*{0.21\paperheight}
    {\color{noirmuted}\small\bfseries SECTION\par}
    \vspace{3mm}
    {\color{white}\bfseries\fontsize{29}{32}\selectfont #1\par}
    \vspace{4mm}
    {\color{noirmuted}\fontsize{13}{16}\selectfont #2\par}
  \end{frame}
}

\newcommand{\eyebrow}[1]{{\color{noircyan}\small\bfseries #1\par}}

\newcommand{\panelbox}[3][panel]{%
  \begin{beamercolorbox}[wd=\linewidth,sep=6pt,rounded=true]{#1}
    {\color{noircyan}\bfseries #2\par}
    \vspace{0.5mm}
    {\small #3\par}
  \end{beamercolorbox}
}

\newcommand{\statcard}[2]{%
  \begin{beamercolorbox}[wd=\linewidth,sep=6pt,rounded=true]{panel alt}
    {\color{noircyan}\bfseries\fontsize{14}{16}\selectfont #1\par}
    \vspace{0.6mm}
    {\color{noirmuted}\scriptsize #2\par}
  \end{beamercolorbox}
}

\newcommand{\blocksep}{%
  \vspace{0.7mm}%
  {\color{noirline}\hrule height 0.35pt}%
  \vspace{0.7mm}%
}

\tikzset{
  flowbox/.style={
    draw=noirline,
    fill=noirpanelalt,
    text=noirtext,
    rounded corners=4pt,
    align=center,
    text width=2.7cm,
    minimum height=0.92cm,
    inner sep=4pt,
    font=\scriptsize
  },
  flowboxwide/.style={
    draw=noirline,
    fill=noirpanelalt,
    text=noirtext,
    rounded corners=4pt,
    align=center,
    text width=3.35cm,
    minimum height=0.92cm,
    inner sep=4pt,
    font=\scriptsize
  },
  repo/.style={
    draw=noircyan!70,
    fill=noirpanel,
    text=noirtext,
    rounded corners=5pt,
    align=left,
    text width=3.55cm,
    minimum height=2.1cm,
    inner sep=7pt,
    font=\scriptsize
  },
  arrow/.style={-Latex, line width=0.8pt, draw=noircyan},
  warnarrow/.style={-Latex, line width=0.8pt, draw=duneorange},
  smallnote/.style={align=center, font=\tiny, text=noirtext}
}

\newcommand{\daqPipelineDiagram}{%
\resizebox{\linewidth}{!}{%
\begin{tikzpicture}[node distance=0.55cm and 0.45cm]
  \node[flowboxwide] (hw) {DAPHNE\\Ethernet output};
  \node[flowboxwide, right=of hw] (daq) {DUNE DAQ\\readout path};
  \node[flowboxwide, right=of daq] (frag) {Fragments\\with DAPHNE ETH frames};
  \node[flowboxwide, right=of frag] (decode) {Python/C++\\decode interface};
  \node[flowboxwide, right=of decode] (analysis) {HDF5/NumPy\\Waffles analysis};
  \draw[arrow] (hw) -- (daq);
  \draw[arrow] (daq) -- (frag);
  \draw[arrow] (frag) -- (decode);
  \draw[arrow] (decode) -- (analysis);
  \node[smallnote, below=0.25cm of daq] {data taking baseline};
  \node[smallnote, below=0.25cm of decode] {reviewer-visible contract};
\end{tikzpicture}
}%
}

\newcommand{\standaloneDiagram}{%
\resizebox{0.9\linewidth}{!}{%
\begin{tikzpicture}[node distance=0.55cm and 0.55cm]
  \node[flowbox] (scripts) {\texttt{daphne-server}\\scan scripts};
  \node[flowbox, right=of scripts] (tunnel) {SSH tunnel\\or direct host};
  \node[flowbox, right=of tunnel] (server) {DAPHNE control\\server};
  \node[flowbox, right=of server] (spy) {Spy buffer\\or osc capture};
  \node[flowbox, below=0.85cm of server] (ssp) {SSP LED\\control};
  \node[flowbox, right=of spy] (products) {run folders\\plots, JSON};
  \draw[arrow] (scripts) -- (tunnel);
  \draw[arrow] (tunnel) -- (server);
  \draw[arrow] (server) -- (spy);
  \draw[arrow] (spy) -- (products);
  \draw[warnarrow] (scripts) |- (ssp);
  \draw[warnarrow] (ssp) -| (spy);
  \node[smallnote, above=0.22cm of server] {no DAQ run control};
  \node[smallnote, below=0.22cm of products] {calibration evidence};
\end{tikzpicture}
}%
}

\newcommand{\pdsRunDiagram}{%
\resizebox{0.94\linewidth}{!}{%
\begin{tikzpicture}[node distance=0.48cm and 0.42cm]
  \node[flowbox] (paths) {\texttt{00\_paths}\\facility};
  \node[flowbox, right=of paths] (defaults) {\texttt{02}/\texttt{03}\\defaults};
  \node[flowbox, right=of defaults] (scan) {scan JSON\\one variable};
  \node[flowbox, right=of scan] (patch) {minimal DAPHNE\\JSON patch};
  \node[flowbox, right=of patch] (xml) {\texttt{add\_daphne\_conf}\\OKS XML};
  \node[flowbox, right=of xml] (ssp) {\texttt{set\_ssp\_conf}\\LED};
  \node[flowbox, right=of ssp] (drunc) {\texttt{drunc}\\short run};
  \draw[arrow] (paths) -- (defaults);
  \draw[arrow] (defaults) -- (scan);
  \draw[arrow] (scan) -- (patch);
  \draw[arrow] (patch) -- (xml);
  \draw[arrow] (xml) -- (ssp);
  \draw[arrow] (ssp) -- (drunc);
  \node[smallnote, below=0.2cm of xml] {overlay, not hand edit};
  \node[smallnote, below=0.2cm of drunc] {data + monitoring};
\end{tikzpicture}
}%
}

\begin{document}

{
\setbeamertemplate{headline}{}
\setbeamertemplate{footline}{}
\begin{frame}[plain,noframenumbering]
  \begin{tikzpicture}[remember picture,overlay]
    \fill[noirbg] (current page.south west) rectangle (current page.north east);
    \fill[noirpanelalt] (current page.north west) rectangle ($(current page.north east)+(0,-2.1)$);
    \fill[duneorange] ($(current page.north west)+(0.8,-0.62)$) rectangle ++(1.5,0.08);
    \fill[noircyan] ($(current page.south east)+(-2.8,0.8)$) rectangle ++(1.9,0.08);
    \IfFileExists{logos/DUNElogo_white.png}{%
      \node[anchor=south east,opacity=0.13] at ($(current page.south east)+(-0.55,1.48)$) {%
        \includegraphics[width=0.33\paperwidth]{logos/DUNElogo_white.png}%
      };
    }{}%
    \IfFileExists{logos/FNAL-Logo-NAL-Light.png}{%
      \node[anchor=south east] at ($(current page.south east)+(-0.92,0.82)$) {%
        \includegraphics[height=0.45cm]{logos/FNAL-Logo-NAL-Light.png}%
      };
    }{}%
  \end{tikzpicture}
  \vspace*{0.26\paperheight}
  {\color{white}\bfseries\fontsize{29}{31}\selectfont DAPHNE DAQ Integration\par}
  \vspace{2mm}
  {\color{noirmuted}\fontsize{15}{18}\selectfont Ethernet data taking, standalone calibration, and \texttt{pds-run} scans\par}
  \vspace{7mm}
  {\color{noircyan}\bfseries Manuel Arroyave\par}
  \vspace{1mm}
  {\color{noirmuted}May 8, 2026\par}
\end{frame}
}

\begin{frame}[t]{Executive Claim}
\begin{columns}[T,onlytextwidth]
\column{0.57\textwidth}
\eyebrow{WHAT IS NOW TRUE}
{\color{white}\bfseries\fontsize{16}{19}\selectfont
DAPHNE Ethernet data taking is established in DAQ; calibration scans are being moved from standalone tools into \texttt{pds-run}. \par}

\vspace{1.5mm}
{\color{noirmuted}\small
Decoding belongs in DAQ libraries, with Waffles downstream. Calibration policy belongs in \texttt{pds-run}. Firmware/schema work adds peak descriptors and the x-correlation self trigger; firmware dead-time remains active. \par}

\column{0.39\textwidth}
\begin{beamercolorbox}[wd=\linewidth,sep=6pt,rounded=true]{panel alt}
\scriptsize
{\color{noircyan}\bfseries Nov 2025}\par
DAQ baseline proven in the VD Coldbox.\par
\blocksep
{\color{noircyan}\bfseries Feb 2026}\par
Standalone calibration proven in the same VD Coldbox facility.\par
\blocksep
{\color{noircyan}\bfseries Mar--Apr 2026}\par
Peak descriptors and x-correlation self trigger implemented; dead-time work ongoing.\par
\blocksep
{\color{noircyan}\bfseries May 2026}\par
DAQ-wrapped calibration exercised with \texttt{pds-run}.\par
\end{beamercolorbox}
\end{columns}
\end{frame}

\begin{frame}[t]{Scope And Success Criteria}
\begin{columns}[T,onlytextwidth]
\column{0.31\textwidth}
\panelbox{DAQ baseline}{%
Proven in the VD Coldbox, November 2025; decodes DAPHNE Ethernet fragments and preserves monitoring in fddaq-v5.5.0.}

\column{0.31\textwidth}
\panelbox{Standalone calibration}{%
Proven in the same VD Coldbox facility, February 2026, with direct DAPHNE control and spy-buffer capture.}

\column{0.31\textwidth}
\panelbox{DAQ-wrapped calibration}{%
May 2026 work: selected scan variables move into \texttt{pds-run}; every point creates a DAQ run.}
\end{columns}

\vspace{2mm}
\begin{beamercolorbox}[wd=\linewidth,sep=6pt,rounded=true]{panel alt}
\footnotesize
\textcolor{noircyan}{\textbf{Constraint:}} the DAQ environment, OKS database, proxy state, and tmux session are part of the apparatus. The change must survive a short run on \texttt{np04-srv-024} through \texttt{np02-vdcb}.
\end{beamercolorbox}
\end{frame}

\sectioncard{DAQ Baseline}{The three DAQ packages define the format, the decode API, and the live controller path}

\begin{frame}[t]{Baseline Data Path}
\centering
\vspace{2mm}
\daqPipelineDiagram

\vspace{3mm}
\begin{columns}[T,onlytextwidth]
\column{0.32\textwidth}
\panelbox{Format}{C++ classes describe the wire layout and enforce bit-field boundaries.}
\column{0.32\textwidth}
\panelbox{Decode}{Python bindings expose arrays for analysis without rewriting the unpacker.}
\column{0.32\textwidth}
\panelbox{Waffles}{Waffles consumes the decoded arrays for calibration and validation workflows.}
\end{columns}
\end{frame}

\begin{frame}[t]{Changed DAQ Repositories}
\begin{columns}[T,onlytextwidth]
\column{0.32\textwidth}
\begin{tikzpicture}
\node[repo] {
{\color{noircyan}\bfseries fddetdataformats}\\[1mm]
\texttt{marroyav/daphneeth}\\
Frame definitions and tests\\[1mm]
\textcolor{duneorange}{New: peak descriptor support}\\
\textcolor{noirmuted}{Adds DAPHNE ETH stream frame bindings.}
};
\end{tikzpicture}

\column{0.32\textwidth}
\begin{tikzpicture}
\node[repo] {
{\color{noircyan}\bfseries rawdatautils}\\[1mm]
\texttt{marroyav/daphneeth}\\
NumPy-facing unpackers\\[1mm]
\textcolor{noirmuted}{Exposes ADC, timestamp, channel, stream, and peak descriptor arrays for Waffles and analysis.}
};
\end{tikzpicture}

\column{0.32\textwidth}
\begin{tikzpicture}
\node[repo] {
{\color{noircyan}\bfseries daphnemodules}\\[1mm]
\texttt{patch/fddaq-v5.5.x}\\
DAPHNE V3 controller\\[1mm]
\textcolor{noirmuted}{Adds V3 GeneralInfo monitoring used during operations.}
};
\end{tikzpicture}
\end{columns}
\end{frame}

\begin{frame}[t]{Data-Format Contract}
\begin{columns}[T,onlytextwidth]
\column{0.5\textwidth}
\begin{beamercolorbox}[wd=\linewidth,sep=4pt,rounded=true]{panel}
\scriptsize
{\color{noircyan}\bfseries\texttt{DAPHNEEthFrame}\par}
Single-channel frame: Ethernet header, DAPHNE metadata, five peak descriptors, and 1024 packed 14-bit ADC samples.
\end{beamercolorbox}

\vspace{1mm}
\begin{beamercolorbox}[wd=\linewidth,sep=4pt,rounded=true]{panel alt}
\scriptsize
{\color{noircyan}\bfseries Peak descriptors\par}
New firmware information is now persisted in the frame and exposed through DAQ formats and Python decode utilities.
\end{beamercolorbox}

\vspace{1mm}
\begin{beamercolorbox}[wd=\linewidth,sep=4pt,rounded=true]{panel}
\scriptsize
{\color{noircyan}\bfseries\texttt{DAPHNEEthStreamFrame}\par}
Four-channel stream frame with 280 samples per channel packed as 14-bit ADCs.
\end{beamercolorbox}

\column{0.46\textwidth}
\begin{beamercolorbox}[wd=\linewidth,sep=4pt,rounded=true]{panel alt}
\scriptsize
\textcolor{noircyan}{\textbf{Reviewer handles}}
\vspace{0.5mm}
\begin{itemize}
  \item Constants: ADC width, sample count, channel count.
  \item Accessors validate channel, sample, and peak indices.
  \item Unit tests round-trip ADC and peak fields.
  \item Bindings register triggered and stream types.
\end{itemize}
\end{beamercolorbox}
\end{columns}
\end{frame}

\begin{frame}[t]{Decode Contract}
\begin{columns}[T,onlytextwidth]
\column{0.46\textwidth}
\eyebrow{RAW DATA UTILS}
\footnotesize
\texttt{rawdatautils} is not a second format definition. It reuses \texttt{fddetdataformats} and exposes arrays consumed by Waffles and validation scripts:

\begin{itemize}
  \item ADC, timestamp, and channel arrays
  \item stream-frame equivalents
  \item peak descriptor arrays for triggered frames
\end{itemize}

\column{0.5\textwidth}
\begin{tikzpicture}[node distance=0.3cm]
  \node[flowboxwide] (frag) {DAQ fragment};
  \node[flowboxwide, below=of frag] (reinterpret) {reinterpret as\\DAPHNE frame};
  \node[flowboxwide, below=of reinterpret] (arrays) {NumPy arrays\\ADC, time, channel};
  \node[flowboxwide, below=of arrays] (analysis) {Waffles\\analysis and validation};
  \draw[arrow] (frag) -- (reinterpret);
  \draw[arrow] (reinterpret) -- (arrays);
  \draw[arrow] (arrays) -- (analysis);
\end{tikzpicture}
\end{columns}
\end{frame}

\begin{frame}[t]{Controller And Monitoring Contract}
\begin{columns}[T,onlytextwidth]
\column{0.5\textwidth}
\panelbox{\texttt{daphnemodules} V3 controller}{%
Configures the DAPHNE V3 analog chain, reads counters, and publishes opmon in DAQ.}

\vspace{1mm}
\panelbox[panel alt]{Important release finding}{%
\texttt{fddaq-v5.5.0-a9-1} ships 3.0.1. V3 \texttt{GeneralInfo} opmon appears in 3.0.2 on \texttt{patch/fddaq-v5.5.x}.}

\column{0.46\textwidth}
\begin{tikzpicture}[node distance=0.3cm]
  \node[flowboxwide] (counter) {ReadTriggerCounters};
  \node[flowboxwide, below=of counter] (perchan) {per-channel\\threshold, record, busy, full};
  \node[flowboxwide, below=of perchan] (general) {ReadGeneralInfo\\bias, power, temp};
  \node[flowboxwide, below=of general] (opmon) {Kafka opmon\\shifter view};
  \draw[arrow] (counter) -- (perchan);
  \draw[arrow] (perchan) -- (general);
  \draw[arrow] (general) -- (opmon);
\end{tikzpicture}
\end{columns}
\end{frame}

\sectioncard{Standalone Calibration}{\texttt{daphne-server} established fast board-level scans outside DAQ run control}

\begin{frame}[t]{Why Standalone Was Useful}
\centering
\vspace{0mm}
\standaloneDiagram

\vspace{1mm}
\begin{columns}[T,onlytextwidth]
\column{0.32\textwidth}
\panelbox{Scientific role}{%
Find settings before DAQ run-control time.}
\column{0.32\textwidth}
\panelbox{Engineering role}{%
Exercise routes, spy buffers, LED clients, and analysis.}
\column{0.32\textwidth}
\begin{beamercolorbox}[wd=\linewidth,sep=6pt,rounded=true]{panel}
\footnotesize
\textcolor{noircyan}{\textbf{Boundary}}\par
\vspace{0.5mm}
Diagnostic data is not a substitute for DAQ-integrated data taking.
\end{beamercolorbox}
\end{columns}
\end{frame}

\begin{frame}[t]{Standalone Products}
\begin{columns}[T,onlytextwidth]
\column{0.33\textwidth}
\statcard{Dark scans}{Scan \texttt{vgain}/offset and find pedestal regions from spy-buffer captures.}

\column{0.33\textwidth}
\statcard{LED scans}{Scan LED intensity, mask, and channel sets with direct SSP control.}

\column{0.33\textwidth}
\statcard{Bundles}{Generate per-\texttt{vgain} JSON settings for later replay.}
\end{columns}

\vspace{3mm}
\begin{beamercolorbox}[wd=\linewidth,sep=6pt,rounded=true]{panel}
\footnotesize
Example dynamic-range result: \texttt{vgain} 1500--2500 is useful for the tested LED response; 2750 and 3000 did not reach the 2000 ADC pedestal target in the scanned offset range.
\end{beamercolorbox}
\end{frame}

\begin{frame}[fragile,t]{Standalone Command Shape}
\small
\begin{lstlisting}[style=terminal]
./utils/run_client_via_tunnel.sh daphne-13 \
  client/configure_fe_from_output.py \
  client/dynamic_range_led/per_channel_target2k_vgain2000.json

./utils/run_client_via_tunnel.sh daphne-13 \
  client/protobuf_acquire_list_channels.py --osc_mode \
  -foldername client/dynamic_range_led/runs/vgain2000_led_sat \
  -channel_list 4,5,6,7,12,13,14,15,18,19,21 -N 64 -L 2048
\end{lstlisting}

\vspace{2mm}
\begin{beamercolorbox}[wd=\linewidth,sep=6pt,rounded=true]{panel alt}
\footnotesize
Operational caution: the VD Coldbox target must be checked in configuration. Inventory maps \texttt{daphne-14} to \texttt{10.73.137.160}; some historical examples use \texttt{daphne-13}/\texttt{10.73.137.161}.
\end{beamercolorbox}
\end{frame}

\sectioncard{DAQ-Wrapped Calibration}{\texttt{pds-run} moves selected scan variables into the DAQ configuration and run-control chain}

\begin{frame}[t]{What \texttt{pds-run} Adds}
\centering
\vspace{0mm}
\hspace*{0.03\linewidth}\pdsRunDiagram

\vspace{1mm}
\begin{columns}[T,onlytextwidth]
\column{0.31\textwidth}
\panelbox{Model}{Paths, commands, defaults, and scan dimensions live in numbered JSON files.}
\column{0.31\textwidth}
\panelbox{Overlay}{DAPHNE changes become minimal JSON patches, then OKS XML.}
\column{0.31\textwidth}
\panelbox{Execution}{Each scan point can set SSP LED state and take a DAQ run.}
\end{columns}
\end{frame}

\begin{frame}[t]{Configuration Chain}
\begin{columns}[T,onlytextwidth]
\column{0.47\textwidth}
\begin{beamercolorbox}[wd=\linewidth,sep=6pt,rounded=true]{panel}
\footnotesize
\textcolor{noircyan}{\textbf{Facility}}
\begin{itemize}
  \item \texttt{00\_paths.json}: workarea, DB, OKS, target.
  \item \texttt{01\_commands.json}: proxy and helper commands.
  \item \texttt{02\_run\_defaults.json}: rate, wait, SSP LED.
  \item \texttt{03\_default\_config.json}: base DAPHNE detail.
\end{itemize}
\end{beamercolorbox}

\column{0.49\textwidth}
\begin{beamercolorbox}[wd=\linewidth,sep=6pt,rounded=true]{panel alt}
\footnotesize
\textcolor{noircyan}{\textbf{Scan}}
\begin{itemize}
  \item \texttt{thr-scan}, \texttt{selfthr-scan}, \texttt{att-scan}.
  \item \texttt{offset-scan}, \texttt{trim-scan}.
  \item \texttt{led-intensity-scan}: masks and intensities.
  \item \texttt{plan\_only}: inspect XML and commands.
\end{itemize}
\end{beamercolorbox}
\end{columns}

\vspace{2mm}
\begin{beamercolorbox}[wd=\linewidth,sep=6pt,rounded=true]{panel}
\footnotesize
Review rule: one scan file explains the changed variable; it does not re-declare the whole DAQ environment.
\end{beamercolorbox}
\end{frame}

\begin{frame}[fragile,t]{DAQ Run Command Generated By \texttt{pds-run}}
\begin{lstlisting}[style=terminal]
drunc-unified-shell ssh-CERN-kafka.json \
  /nfs/sw/dunedaq/dunedaq-fddaq-v5.5.0-dev-pds/.../pds-vdcb-session.data.xml \
  np02-session pds-vdcb-pds \
  start-run --run-type TEST \
  change-rate --trigger-rate 20.0 \
  wait 120 shutdown terminate
\end{lstlisting}

\vspace{2mm}
\begin{columns}[T,onlytextwidth]
\column{0.32\textwidth}
\panelbox{Before run}{Source DAQ env, source DB path, disable proxy for runtime.}
\column{0.32\textwidth}
\panelbox{During run}{Run under tmux on the run-control host.}
\column{0.32\textwidth}
\panelbox{After run}{Confirm data files and opmon.}
\end{columns}
\end{frame}

\begin{frame}[fragile,t]{Operations Commands For \texttt{pds-run} Scans}
\begin{lstlisting}[style=terminal]
# Plan first: set "plan_only": true in the scan JSON.
pds-run led-intensity-scan configs/vd_coldbox/06_led_calib.json
pds-run led-intensity-scan configs/vd_coldbox/02_run_defaults.json \
  --mask-values 1,2,4,8 \
  --led-intensity-values 3000,3500,4095
pds-run thr-scan configs/vd_coldbox/04_stthre.json
pds-run selfthr-scan configs/vd_coldbox/04_stthre.json
pds-run offset-scan configs/vd_coldbox/<offset_scan>.json
pds-run trim-scan configs/vd_coldbox/<trim_scan>.json
\end{lstlisting}

\vspace{0.5mm}
\begin{beamercolorbox}[wd=\linewidth,sep=5pt,rounded=true]{panel alt}
\footnotesize
Each command updates only the scan variable, applies the minimal DAPHNE XML overlay, runs SSP setup when requested, then delegates to \texttt{drunc}.
\end{beamercolorbox}
\end{frame}

\begin{frame}[t]{Operational Environment}
\begin{columns}[T,onlytextwidth]
\column{0.48\textwidth}
\panelbox{Login path}{%
\texttt{marroyav@np04-srv-017} then \texttt{np04daq@np04-srv-024}; use tmux session \texttt{np02-vdcb}.}
\vspace{1mm}
\panelbox{Workarea}{%
\texttt{/nfs/sw/dunedaq/}\\
\texttt{dunedaq-fddaq-v5.5.0-dev-pds}.\\
Source \texttt{setup\_for\_run.sh}.}

\column{0.48\textwidth}
\panelbox{Proxy state}{%
Use \texttt{source \textasciitilde/bin/web\_proxy.sh} only for package downloads. Disable for running with \texttt{-u}.}
\vspace{1mm}
\panelbox{Hardware target}{%
Board identity and IP must be checked in the active JSON/XML; target DAPHNE selection is a configuration requirement, not an assumption.}
\end{columns}
\end{frame}

\sectioncard{Validation And Review}{The success condition is both software cleanliness and live DAQ behavior}

\begin{frame}[t]{Live Validation Result}
\begin{columns}[T,onlytextwidth]
\column{0.24\textwidth}
\statcard{Run 44355}{Default release with no local \texttt{daphnemodules}: data taking worked, but V3 monitoring was not visible.}
\column{0.24\textwidth}
\statcard{Run 44356}{Clean \texttt{daphnemodules} 3.0.2 override: two-minute TEST run succeeded at 20 Hz.}
\column{0.24\textwidth}
\statcard{Firmware}{Peak descriptors were added in firmware and then carried into the DAQ data-format contract.}
\column{0.24\textwidth}
\statcard{Conclusion}{Monitoring loss was traced to the release package version, not to run control.}
\end{columns}

\vspace{3mm}
\begin{beamercolorbox}[wd=\linewidth,sep=6pt,rounded=true]{panel alt}
\footnotesize
Final tested state: \texttt{sourcecode/daphnemodules} commit \texttt{13289b5} on \texttt{patch/fddaq-v5.5.x}; \texttt{DAPHNE\_V3\_ROUTE} unset; monitoring visible.
\end{beamercolorbox}
\end{frame}

\begin{frame}[t]{Live \texttt{pds-run} Scan Exercise}
\begin{beamercolorbox}[wd=\linewidth,sep=3pt,rounded=true]{panel}
\tiny
\setlength{\tabcolsep}{3pt}
\renewcommand{\arraystretch}{0.98}
\begin{tabular}{@{}p{0.18\linewidth}p{0.19\linewidth}p{0.54\linewidth}@{}}
\toprule
\textbf{Mode} & \textbf{Runs} & \textbf{Variable exercised} \\
\midrule
self-trigger & 44357--44358 & DAPHNE threshold: 11, 10 \\
AFE attenuation & 44359--44360 & all AFE attenuators: 2200, 2300 \\
offset & 44361--44362 & selected channel offsets: 2260, 2275 \\
trim & 44363--44364 & selected channel trims: 1, 0 \\
AFE bias & 44365--44366 & AFE0 bias: 10, 0; others fixed at 0 \\
LED / SSP & 44367--44368 & LED intensity: 3000, 4095; mask 4 \\
\bottomrule
\end{tabular}
\end{beamercolorbox}

\vspace{1mm}
{\color{noirmuted}\scriptsize
Branch/host: \texttt{marroyav/vdcb-daq-scans} at \texttt{06efc5c} on \texttt{np04-srv-024}, tmux \texttt{np02-vdcb}; logs \texttt{/tmp/pds\_vdcb\_live\_20260508\_152420}.\\[-0.2mm]
Result: every successful scan command exited with status 0; \texttt{xmllint --noout} passed after each XML update; no successful-run traceback or drunc command failure was found.\par}
\end{frame}

\begin{frame}[t]{Live Exercise Findings}
\begin{beamercolorbox}[wd=\linewidth,sep=4pt,rounded=true]{panel}
\scriptsize
\renewcommand{\arraystretch}{1.16}
\begin{tabular}{@{}>{\raggedright\arraybackslash}p{0.25\linewidth}>{\raggedright\arraybackslash}p{0.65\linewidth}@{}}
\toprule
\textbf{Finding} & \textbf{Reviewer-facing consequence} \\
\midrule
SSP object name & Active XML uses \texttt{np02-ssp}; scan configs must not silently rely on the old \texttt{np02-ssp-on} default. \\
Execution seed & \texttt{99\_apply\_fe\_min\_v2\_subset.json} is plan-safe but not execution-safe; \texttt{add\_daphne\_conf} requires \texttt{detector\_id}. \\
Zero/default XML & \texttt{trim=0}, \texttt{v\_bias=0}, and \texttt{bias\_ctrl=0} can be omitted by \texttt{add\_daphne\_conf}; absence of these attributes is not by itself a run failure. \\
Post-scan state & The exercise left SSP at \texttt{channel\_mask=4}, \texttt{pulse1\_width=3}; scan plans should state whether to restore the baseline. \\
\bottomrule
\end{tabular}
\end{beamercolorbox}

\vspace{1.5mm}
\begin{beamercolorbox}[wd=\linewidth,sep=4pt,rounded=true]{panel alt}
\scriptsize
Conclusion: the DAQ-wrapped scan loop is operational for the tested modes; remaining cleanup is configuration hygiene, explicit target naming, and post-scan restoration policy.
\end{beamercolorbox}
\end{frame}

\begin{frame}[plain]
\begin{tikzpicture}[remember picture,overlay]
  \fill[noirpanelalt] (current page.south west) rectangle (current page.north east);
  \fill[duneorange] ($(current page.north west)+(0.9,-0.9)$) rectangle ++(1.8,0.08);
  \fill[noircyan] ($(current page.south east)+(-2.8,0.9)$) rectangle ++(1.8,0.08);
\end{tikzpicture}
\vspace*{0.22\paperheight}
{\color{white}\bfseries\fontsize{30}{34}\selectfont Summary\par}
\vspace{5mm}
{\color{noirmuted}\Large
DAQ baseline: proven in the VD Coldbox, November 2025.\\[2mm]
Standalone calibration: proven in the same facility, February 2026.\\[2mm]
\texttt{pds-run}: DAQ-wrapped calibration, May 2026.\\[2mm]
Peak descriptors: firmware, DAQ format, decode, Waffles.
\par}
\end{frame}

\newcommand{\backupresourcebar}[4]{%
  \begin{tikzpicture}[x=0.1cm,y=0.34cm,baseline=(label.base)]
    \pgfmathsetmacro{\barend}{#2/10}
    \node[anchor=east,font=\scriptsize,text=noirmuted] (label) at (0,0) {#1};
    \draw[draw=noirline,fill=noirpanelalt,rounded corners=1pt] (0.25,-0.17) rectangle (10.25,0.17);
    \draw[draw=#4,fill=#4!30,rounded corners=1pt] (0.25,-0.17) rectangle ({0.25+\barend},0.17);
    \node[anchor=west,font=\scriptsize,text=noirtext] at (10.45,0) {#3};
  \end{tikzpicture}\par
}

\begin{frame}[noframenumbering,t]{Backup: Interface Reference}
\begin{columns}[T,onlytextwidth]
\column{0.48\textwidth}
\eyebrow{FORMAL INTERFACE HANDLE}
{\color{white}\bfseries DUNE Interface Control Document: DAQ \& Photon Detector System (PDS)\par}
\vspace{1mm}
{\color{noirmuted}\small
EDMS 2088726 v7\\
\href{https://edms.cern.ch/document/2088726/7}{https://edms.cern.ch/document/2088726/7}\par}

\column{0.48\textwidth}
\begin{beamercolorbox}[wd=\linewidth,sep=6pt,rounded=true]{panel alt}
\footnotesize
\textcolor{noircyan}{\textbf{Use in discussion}}\par
\vspace{0.5mm}
The ICD is the external reference for the DAQ/PDS boundary. Resource and dead-time work remains an implementation risk, but the mitigation path is constrained by this interface contract.
\end{beamercolorbox}
\end{columns}

\vspace{2mm}
\begin{beamercolorbox}[wd=\linewidth,sep=5pt,rounded=true]{panel}
\scriptsize
Backup source for resource evidence: \texttt{repo/Latex/daphne\_hardware\_architecture\_may2026}, sections ``Baseline PL Resource Report'' and ``Current Grouped Candidate Build Result''.
\end{beamercolorbox}
\end{frame}

\begin{frame}[noframenumbering,t]{Backup: K26 Resource Mitigation Evidence}
\begin{columns}[T,onlytextwidth]
\column{0.49\textwidth}
\eyebrow{BASELINE, OUTPUT-2C337B5}
\vspace{1mm}
\backupresourcebar{CLB sites}{99.55}{99.55\%}{duneorange}
\backupresourcebar{DSP48E2}{99.36}{99.36\%}{duneorange}
\backupresourcebar{CLB LUTs}{74.48}{74.48\%}{noircyan}
\backupresourcebar{Block RAM}{68.75}{68.75\%}{noircyan}
\backupresourcebar{URAM}{62.50}{62.50\%}{noircyan}
\backupresourcebar{Registers}{44.78}{44.78\%}{noircyan}

\column{0.49\textwidth}
\eyebrow{GROUPED CANDIDATE, E5E2D70}
\vspace{1mm}
\backupresourcebar{CLB LUTs}{76.05}{76.05\%}{noircyan}
\backupresourcebar{Block RAM}{67.36}{67.36\%}{noircyan}
\backupresourcebar{RAMB36/FIFO}{63.89}{63.89\%}{noircyan}
\backupresourcebar{URAM}{62.50}{62.50\%}{noircyan}
\backupresourcebar{Registers}{61.19}{61.19\%}{noircyan}
\backupresourcebar{DSP48E2}{0.00}{0.00\%}{noircyan}
\end{columns}

\vspace{2mm}
\begin{beamercolorbox}[wd=\linewidth,sep=5pt,rounded=true]{panel alt}
\scriptsize
Interpretation: the grouped candidate keeps LUT/BRAM/URAM below device limits and removes the baseline DSP pressure. Remaining qualification is Hermes readout, backpressure, and real-trigger behavior, not FPGA fit or route legality.
\end{beamercolorbox}
\end{frame}

\end{document}
